\documentclass[12pt,a4paper]{article}
\usepackage[margin=2cm]{geometry}
\usepackage{amsmath,amssymb}
\usepackage{fontspec}
\setmainfont{DejaVu Serif}
\setlength{\parindent}{0pt}
\setlength{\parskip}{4pt}
\begin{document}
\section*{Giải và biện luận hệ phương trình}
% line:l1
Giải và biện luận hệ phương trình sau theo tham số m ∈ ℝ:

% line:l2
\[ \begin{cases}x+y+2z=3\\2x+3y+5z=7\\x+y+\left(m^{2}+1\right)z=m+2\end{cases} \]

% line:l3
Ma trận bổ sung của hệ:

% line:l4
\[ \overline{A}=\left[\begin{array}{ccc|c}1&1&2&3\\2&3&5&7\\1&1&m^{2}+1&m+2\end{array}\right] \]

% line:l5
\[ \xrightarrow[d_{3}\to d_{3}-d_{1}]{d_{2}\to d_{2}-2d_{1}}\left[\begin{array}{ccc|c}1&1&2&3\\0&1&1&1\\0&0&m^{2}-1&m-1\end{array}\right] \]

% line:l6
Biện luận:

\begin{align*}
% line:l7
& m^{2}-1=0 \\
% line:l8
& \Leftrightarrow m=\pm1
\end{align*}

% line:l9
• Nếu m ≠ ±1 thì hệ có nghiệm duy nhất:

% line:l10
\[ z=\frac{m-1}{m^{2}-1}=\frac{1}{m+1},\quad y=1-z=\frac{m}{m+1},\quad x=3-y-2z=\frac{2m+1}{m+1} \]

% line:l11
• Nếu m = 1: hệ có vô số nghiệm

% line:l12
\[ \begin{cases}x=2-z\\y=1-z\\z\in\mathbb{R}\end{cases} \]

% line:l13
• Nếu m = −1: phương trình cuối trở thành 0z = −2, hệ vô nghiệm.

\end{document}
